\chapter{Implementation}
\label{ch:implementation}

This chapter presents the concrete implementation of the system
described in Chapter~\ref{ch:methodology}. It is organised top-down:
we first describe the project's directory structure so that the
reader can locate any referenced file, then walk through the key
modules in three programming languages (Python for the ML pipeline,
Kotlin for the Wear~OS application, and C++ for the Arduino firmware),
and finally describe the simulation framework that allows the entire
end-to-end system to be exercised without physical hardware.

\section{Project Structure}
\label{sec:impl-structure}

The project is organised into four top-level directories that
correspond to the four major build artefacts (Python wheel, Wear~OS
APK, Arduino sketch, and simulation scripts):

\begin{lstlisting}[language=bash]
stress_detection/               # Python 3.11 ML pipeline
  data/
    wesad_loader.py             # WESAD pickle I/O
    feature_engineering.py      # resampling + IBI + label mapping
    dataset_builder.py          # windowing + train/test split
  model/
    stress_model.py             # BiLSTM + attention definition
    attention_layer.py          # custom Keras attention layer
  training/
    train.py                    # end-to-end training script
    evaluate.py                 # held-out evaluation + confusion matrix
  conversion/
    export_tflite.py            # SavedModel -> TFLite conversion
  models/
    stress_model.tflite         # deployable FlatBuffer (72 KB)
wear_os_app/
  app/src/main/java/com/alcowatch/stress/
    StressInferenceEngine.kt    # TFLite interpreter wrapper
    SensorCollector.kt          # 4 Hz windowing pipeline
    StressBLECharacteristic.kt  # GATT peripheral
    StressLevel.kt              # ordinal enum
arduino/
  stress_cabin_control/
    stress_cabin_control.ino    # top-level sketch
    cabin_profile.h             # profile constants and LED tables
    ble_central.cpp             # ArduinoBLE central wrapper
scripts/
  run_simulation.py             # scripted BLE central for hardware-free tests
  generate_figures.py           # paper figure generation
paper_figures/                  # output PDFs
\end{lstlisting}

The separation of \texttt{stress\_detection/} from the
\texttt{wear\_os\_app/} and \texttt{arduino/} trees reflects the
build-time boundary between training (performed once, on a
workstation) and deployment (performed continuously, on the device).
The \texttt{models/} subdirectory is the handoff point: its
\texttt{stress\_model.tflite} artefact is the only file that crosses
from training to deployment.

\section{ML Pipeline Implementation}
\label{sec:impl-ml}

The ML pipeline is implemented in Python 3.11 against TensorFlow
2.15. We describe three representative modules: the WESAD loader,
the feature-engineering module that performs label mapping, and the
model definition itself.

\subsection{\texttt{wesad\_loader.py}}

WESAD is distributed as per-subject pickle files containing a nested
dictionary of signals and labels. Python 3 cannot load these pickles
with the default encoding because they were serialised under Python
2; an explicit \texttt{encoding='latin1'} argument is required.

\begin{lstlisting}[language=Python]
import pickle
from pathlib import Path
from typing import Dict

class WESADLoader:
    """Load per-subject WESAD pickle files."""

    def __init__(self, root: Path):
        self.root = Path(root)

    def load_subject(self, subject_id: str) -> Dict:
        """Load one subject (e.g. 'S2') from the WESAD release.

        Returns a dict with keys 'signal' (nested wrist/chest) and
        'label' (1D ndarray at 700 Hz).
        """
        path = self.root / subject_id / f"{subject_id}.pkl"
        with open(path, "rb") as fh:
            # WESAD was pickled under Python 2; latin1 decoding is
            # required to restore the byte strings as expected.
            data = pickle.load(fh, encoding="latin1")
        return {
            "bvp":   data["signal"]["wrist"]["BVP"].ravel(),    # 64 Hz
            "eda":   data["signal"]["wrist"]["EDA"].ravel(),    # 4 Hz
            "temp":  data["signal"]["wrist"]["TEMP"].ravel(),   # 4 Hz
            "acc":   data["signal"]["wrist"]["ACC"],            # 32 Hz, 3-axis
            "label": data["label"].astype("int8").ravel(),       # 700 Hz
        }
\end{lstlisting}

\subsection{\texttt{feature\_engineering.py}}

The label-mapping function implements the stress-split rule from
Section~\ref{sec:m-labels} directly. Per-subject EDA and IBI
quantiles are computed once on the subject's stress-labelled samples
and then applied to each sample in turn.

\begin{lstlisting}[language=Python]
import numpy as np

# WESAD code -> our ordinal stress class (default target)
_BASE_MAP = {1: 0, 4: 0, 3: 1, 2: 2}   # 0 drops; 2 may be bumped to 3

def map_wesad_labels(
    wesad_labels: np.ndarray,
    eda: np.ndarray,
    ibi: np.ndarray,
) -> np.ndarray:
    """Map WESAD 0..4 labels to our 0..3 ordinal stress classes.

    Stress samples (code 2) split into Moderate (2) vs Critical (3)
    by EDA > 75th-pct AND IBI < 25th-pct (per-subject quantiles).
    """
    out = np.full_like(wesad_labels, fill_value=-1, dtype=np.int8)

    stress_mask = wesad_labels == 2
    if stress_mask.any():
        eda_thresh = np.quantile(eda[stress_mask], 0.75)
        ibi_thresh = np.quantile(ibi[stress_mask], 0.25)
        crit_mask = stress_mask & (eda > eda_thresh) & (ibi < ibi_thresh)
    else:
        crit_mask = np.zeros_like(stress_mask)

    for code, cls in _BASE_MAP.items():
        out[wesad_labels == code] = cls
    out[crit_mask] = 3      # promote to Critical
    return out              # -1 entries are later discarded
\end{lstlisting}

\subsection{\texttt{stress\_model.py}}

The model definition is a direct Keras functional-API translation
of Table~\ref{tab:layers}. The attention layer is factored out into
a thin wrapper for readability.

\begin{lstlisting}[language=Python]
import tensorflow as tf
from tensorflow.keras import layers, Model

def build_model(timesteps: int = 30, n_features: int = 5,
                n_classes: int = 4) -> Model:
    """BiLSTM(64) + additive attention + 2-layer MLP head.

    Approx. 160k parameters, 72 KB after TFLite default-optimisation.
    """
    inputs = layers.Input(shape=(timesteps, n_features),
                          name="stress_input")

    # Bidirectional temporal encoder
    x = layers.Bidirectional(
        layers.LSTM(64, return_sequences=True), name="bilstm"
    )(inputs)                                  # -> [B, T, 128]
    x = layers.Dropout(0.3)(x)

    # Additive attention pooling
    e = layers.Dense(1, activation="tanh", name="attn_score")(x)   # [B, T, 1]
    a = layers.Softmax(axis=1, name="attn_weights")(e)             # [B, T, 1]
    context = layers.Lambda(
        lambda z: tf.reduce_sum(z[0] * z[1], axis=1),
        name="attn_context"
    )([x, a])                                                      # [B, 128]

    # Feed-forward head
    h = layers.Dense(64, activation="relu")(context)
    h = layers.Dropout(0.3)(h)
    h = layers.Dense(32, activation="relu")(h)
    logits = layers.Dense(n_classes, activation="softmax",
                          name="stress_probs")(h)

    model = Model(inputs, logits, name="stress_classifier")
    model.compile(
        optimizer=tf.keras.optimizers.Adam(1e-3),
        loss="categorical_crossentropy",
        metrics=["accuracy"],
    )
    return model
\end{lstlisting}

\section{Wear OS Kotlin Implementation}
\label{sec:impl-kotlin}

The Wear~OS application exposes two Kotlin classes directly relevant
to this report: \texttt{StressInferenceEngine}, which wraps the
TensorFlow Lite interpreter, and \texttt{StressBLECharacteristic},
which builds the twelve-byte GATT payload.

\subsection{\texttt{StressInferenceEngine.kt}}

\begin{lstlisting}[language=Kotlin]
class StressInferenceEngine(context: Context) {

    private val interpreter: Interpreter =
        Interpreter(loadModelFile(context, "stress_model.tflite"))

    // Normalisation constants exported from training.
    private val featureMeans  = floatArrayOf( /* 5 values */ )
    private val featureStds   = floatArrayOf( /* 5 values */ )

    /** Run one inference on a 30x5 window. */
    fun classify(window: Array<FloatArray>): StressResult {
        require(window.size == 30 && window[0].size == 5)

        // z-score normalise in place
        val normalised = Array(30) { t ->
            FloatArray(5) { f ->
                (window[t][f] - featureMeans[f]) / featureStds[f]
            }
        }

        // TFLite expects [1, 30, 5] input, [1, 4] output
        val input  = arrayOf(normalised)
        val output = Array(1) { FloatArray(4) }
        interpreter.run(input, output)

        val probs = output[0]
        val level = probs.indices.maxBy { probs[it] }
        return StressResult(
            level       = StressLevel.fromInt(level),
            confidence  = probs[level],
            probabilities = probs
        )
    }
}
\end{lstlisting}

\subsection{\texttt{StressBLECharacteristic.kt}}

\begin{lstlisting}[language=Kotlin]
object StressBLECharacteristic {

    /** Pack a StressResult into the 12-byte GATT payload. */
    fun buildPacket(result: StressResult, timestampMs: Long): ByteArray {
        val buf = ByteBuffer.allocate(12).order(ByteOrder.LITTLE_ENDIAN)
        buf.putLong(timestampMs)                                  // bytes 0..7
        buf.put(result.level.ordinalByte)                          // byte 8
        buf.put((result.confidence * 100f)
                .coerceIn(0f, 100f).toInt().toByte())              // byte 9
        buf.putShort(0)                                            // bytes 10..11 reserved
        return buf.array()
    }
}
\end{lstlisting}

\section{Arduino Firmware}
\label{sec:impl-arduino}

The Arduino firmware is a single-sketch C++ program built against
the \texttt{ArduinoBLE} library. The main loop polls the BLE central,
applies timeout logic, and dispatches to a per-profile actuator
routine.

\subsection{\texttt{loop()}}

\begin{lstlisting}[language=C++]
void loop() {
    BLE.poll();

    if (peripheral && peripheral.connected()) {
        if (stressChar.valueUpdated()) {
            uint8_t buf[12];
            stressChar.readValue(buf, 12);
            // bytes 0..7: timestamp (unused on MCU)
            uint8_t level      = buf[8];
            uint8_t conf_pct   = buf[9];
            applyProfile((StressLevel)level, conf_pct);
            lastPacketMs = millis();
        }
    }

    // 60-second BLE timeout -> WAITING
    if (millis() - lastPacketMs > 60000UL) {
        applyProfile(WAITING, 0);
    }
}
\end{lstlisting}

\subsection{\texttt{applyProfile()}}

\begin{lstlisting}[language=C++]
static void applyProfile(StressLevel s, uint8_t confPct) {
    switch (s) {
        case CALM:     // warm white, 30% fan, audio off
            setRGB(255, 180,  90); setFan( 76); audioOff();         break;
        case MILD:     // neutral white, 50% fan, audio off
            setRGB(255, 255, 220); setFan(128); audioOff();         break;
        case MODERATE: // blue, 75% fan, soft music
            setRGB( 40, 120, 255); setFan(191); audioSoftMusic();   break;
        case CRITICAL: // red blinking, 100% fan, voice alert
            blinkRGB(255,  20,  20, 500); setFan(255); audioVoice();break;
        case WAITING:  // yellow blink, fan off, audio off
            blinkRGB(255, 200,   0, 1000); setFan(  0); audioOff(); break;
    }
    currentProfile = s;
}
\end{lstlisting}

\section{Simulation Framework}
\label{sec:impl-simulation}

A hardware-free simulation framework allows the Arduino-side firmware
and the full BLE round-trip to be exercised without any physical
hardware. \texttt{scripts/run\_simulation.py} implements a scripted
BLE central (via \texttt{bleak}) that connects to a virtual wearable,
emits a pre-recorded twelve-byte packet stream, and asserts on the
Arduino-side serial output. Table~\ref{tab:scenarios} lists the five
integration scenarios exercised.

\begin{table}[H]
\centering
\caption{Simulation scenarios for the stress-detection pipeline. All
five scenarios are run on every CI build.}
\label{tab:scenarios}
\small
\begin{tabular}{clll}
\toprule
\textbf{\#} & \textbf{Name} & \textbf{Trigger} & \textbf{Expected cabin
response} \\
\midrule
7  & Calm commute    & level=0 for 2 min   & warm white, fan 30\%, audio off \\
8  & Mild arousal    & level=1 for 2 min   & neutral white, fan 50\%, audio off \\
9  & Moderate stress & level=2 for 2 min   & blue, fan 75\%, soft music \\
10 & Critical event  & level=3 for 30 s    & red blinking, fan 100\%, voice \\
11 & Watch removed   & BLE disconnect 90 s & yellow blink (WAITING), fan 0\% \\
\bottomrule
\end{tabular}
\end{table}

These five scenarios collectively cover the complete FSM transition
graph: every edge is traversed at least once. Scenario 11 additionally
exercises the sixty-second timeout path, which is the only transition
not driven by an incoming BLE packet. Regression on any scenario causes
the CI build to fail, which enforces continued correctness of both
the Arduino state machine and the Wear~OS BLE packet layout.
